\section{Discussion}
\label{sec:limits}
OMP, Pi-fixture, and REST routes completed the evaluated integration traces
in two simulator scenes. This case study complements asset-specific approval
and process-validation studies \cite{graul2026aas,schall2026safe}, but
measures neither process performance nor comparative platform value.

One simulator and one probe per task preclude independent process
replication. Only reported adapter configurations were tested; general
interchangeability, model behavior, online DAQ, configuration reuse, physical
transfer, and a platform baseline remain untested. The tuple check guards
software dispatch, not device I/O. Synthetic approval does not assess operator
performance \cite{parasuramanmanzey2010complacency,parasuraman2000model,bainbridge1983ironies};
operation restrictions do not establish security
\cite{greshake2023injection,liu2024pinjection,iec62443-3-3} or replace PLC/SIS
safeguards \cite{iec61511}. Future work needs a defined comparison and
independent simulator/plant trials.

Scope and channel-policy refusals are not uniformly persisted; auditability
is limited to the records and denial classes described in Section~IV.

\section{Conclusion}
\label{sec:conclusion}
AgentWorkShop unifies registered-task and direct-call access to shared
industrial services while retaining route-specific records. Tests across two
simulator mappings verify selected traces and one pre-dispatch context guard,
not process improvement or comparative platform benefit.
